% !Mode:: "TeX:UTF-8"
\documentclass{mcmthesis}
%\documentclass[CTeX = true]{mcmthesis}  % 当使用 CTeX 套装时请注释上一行使用该行的设置
\setlength{\headheight}{15pt}
\mcmsetup{tstyle=\color{red}\bfseries,%修改题号，队号的颜色和加粗显示，黑色可以修改为 black
	tcn = 2604335, problem = B, %修改队号，参赛题号
	sheet = true, titleinsheet = true, keywordsinsheet = true,
	titlepage = false, abstract = true}
%四款字体可以选择
%\usepackage{times}
%\usepackage{newtxtext}
%\usepackage{palatino}
\usepackage{txfonts}
\usepackage {booktabs}
\usepackage {caption}
\usepackage{indentfirst}  %首行缩进，注释掉，首行就不再缩进。
\usepackage{lipsum}
\title{LATEX TITLE}
\author{\small \href{https://www.latexstudio.net/}
	{\includegraphics[width=7cm]{mcmthesis-logo}}}
\date{\today}
\begin{document}
	\begin{abstract}
		\par Use this template to begin typing the first page (summary page) of your electronic report. This template uses a 12-point Times New Roman font. Submit your paper as an Adobe PDF electronic file (e.g. 1111111.pdf), typed in English, with a readable font of at least 12-point type.
		
		Do not include the name of your school, advisor, or team members on this or any page.
		
		Papers must be within the 25 page limit.
		
		Be sure to change the control number and problem choice above.
		You may delete these instructions as you begin to type your report here.
		
		Follow us @COMAPMath on Twitter or COMAPCHINAOFFICIAL on Weibo for the most up to date contest information.
		
		\begin{keywords}
			keyword1; keyword2
		\end{keywords}
	\end{abstract}
	\maketitle
	%% Generate the Table of Contents, if it's needed.
	\tableofcontents
	\newpage
	%%
	%% Generate the Memorandum, if it's needed.
	%% \memoto{\LaTeX{}studio}
	%% \memofrom{Liam Huang}
	%% \memosubject{Happy \TeX{}ing!}
	%% \memodate{\today}
	%% \memologo{\LARGE I'm pretending to be a LOGO!}
	%% \begin{memo}[Memorandum]
		%%   \lipsum[1-3]
		%% \end{memo}
	%%
	\section{Introduction}
	\subsection{Problem Background}
    \lipsum[3]
	\subsection{Restatement of the Problem}
	\lipsum[3]
	\subsection{Literature Review}
	\lipsum[3]
	\begin{itemize}
		\item the angular velocity of the bat,
		\item the velocity of the ball, and
		\item the position of impact along the bat.
	\end{itemize}
	\lipsum[4]
	\emph{center of percussion} [Brody 1986], \lipsum[5]
	
	\begin{Theorem} \label{thm:latex}
		\LaTeX
	\end{Theorem}
	\begin{Lemma} \label{thm:tex}
		\TeX .
	\end{Lemma}
	\begin{proof}
		The proof of theorem.
	\end{proof}
	
	\subsection{Our Work}
	\lipsum[6]
	\begin{itemize}
		\item
		\item
		\item
		\item
	\end{itemize}
	
	\lipsum[7]
	
	\section{Assumptions and Justifications}
	\begin{itemize}
	\item Assumption 1: ...
	\item Assumption 2: ...
	\end{itemize}
	\section{Notations}
	The key mathematical notations used in this paper are listed in Table 1.
	\begin {table}[htbp]
	\centering
	\caption {Major notations}
	\begin {tabular}{l c r} % 三列，左、中、右对齐
	\toprule [1pt] % 最上面的粗线
	\textbf {Symbol} & \textbf {Description} & \textbf {Unit} \\
	\midrule % 中间的细线
	x & Example variable & m \\
	v & Velocity & m/s \\
	\bottomrule [1pt] % 最下面的粗线
	\end {tabular}
	\end {table}
	\section{The Preparation of Model}
	\subsection{Data Overview}
	\subsection{Data Collection}
	\section{The Model }
	\lipsum[6]
	\section{Model}
	\lipsum[9]
	\section{Sensitivity Analysis}
	\lipsum[6]
	\section{Expansion of the Model}
	\section{Model Evaluation}
	\subsection{Strengths}
	\lipsum[4]
	\subsection{weaknesses}
	\begin{itemize}
		\item \textbf{Applies widely}\\
		This  system can be used for many types of airplanes, and it also
		solves the interference during  the procedure of the boarding
		airplane,as described above we can get to the  optimization
		boarding time.We also know that all the service is automate.
		\item \textbf{Improve the quality of the airport service}\\
		Balancing the cost of the cost and the benefit, it will bring in
		more convenient  for airport and passengers.It also saves many
		human resources for the airline.
	\end{itemize}
	
	AI cite use \AIcite{AI1,AI2,AI3}
	
	%\section{References}  % article类的section是最高级，自动进目录
	\cleardoublepage 
	\addcontentsline{toc}{section}{References}
	\begin{thebibliography}{99}
		\bibitem{1} D.~E. KNUTH, The \TeX{}book, The American Mathematical Society and Addison-Wesley Publishing Company, 1984-1986.
		\bibitem{2} Lamport, Leslie, \LaTeX{}: ``A Document Preparation System'', Addison-Wesley Publishing Company, 1986.
		\bibitem{3} \url{https://www.latexstudio.net/}
	\end{thebibliography}
	\begin{appendices}
		
		\section{First appendix}
		
		In addition, your report must include a letter to the Chief Financial Officer (CFO) of the Goodgrant Foundation, Mr. Alpha Chiang, that describes the optimal investment strategy, your modeling approach and major results, and a brief discussion of your proposed concept of a return-on-investment (ROI). This letter should be no more than two pages in length.
		
		\begin{letter}{Dear, Mr. Alpha Chiang}
			
			\lipsum[1-2]
			
			\vspace{\parskip}
			
			Sincerely yours,
			
			Your friends
			
		\end{letter}
		Here are simulation programmes we used in our model as follow.\\
		
		\textbf{\textcolor[rgb]{0.98,0.00,0.00}{Input matlab source:}}
		
		\section{Second appendix}
		
		some more text \textcolor[rgb]{0.98,0.00,0.00}{\textbf{Input C++ source:}}
		
	\end{appendices}
	\AImatter
	\begin{ReportAiUse}{9}
		\bibitem{AI1}
		Bing AI\\
		Query1: write a sort algorithm\\
		Output: Sure, I can help you with that. Here is a simple implementation of the selection sort algorithm in Python:
		\begin{lstlisting}[language=python]
			def selection_sort(arr):
			for i in range(len(arr)):
			min_idx = i
			for j in range(i + 1, len(arr)):
			if arr[min_idx] > arr[j]:
			min_idx = j
			arr[i], arr[min_idx] = arr[min_idx], arr[i]
			return arr
		\end{lstlisting}
		This algorithm sorts an array by repeatedly finding the minimum element from the unsorted part of the array and putting it at the beginning. \href{https://www.geeksforgeeks.org/sorting-algorithms/}{The time complexity of this algorithm is O($n^2$)}.
		
		There are many other sorting algorithms available such as bubble sort, insertion sort, merge sort, quick sort, heap sort, and more. Each of these algorithms has its own advantages and disadvantages. You can learn more about them on GeeksforGeeks and Programiz.
		
		I hope this helps!
		
		\bibitem{AI2}
		OpenAl Enie (Nov 5, 2023 version, Ernie 4.0)\\
		Query1: <insert the exact wording of any subsequent input into the Al tool>\\
		Output: <insert the complete output from the second query>
		
		\bibitem{AI3}
		Github CoPilot (Feb 3, 2024 version)\\
		Query1: <insert the exact wording you input into the Al tool>\\
		Output: <insert the complete output fron the Al tool>
		
		\bibitem{AI4}
		Google Bard (Feb 2,2024 version)\\
		Query: <insert the exact wording of your query>\\
		Output: <irsert the complete output from the Al tool>
		
		\bibitem{AI5}
		Bing AI\\
		Query: write a FFT algorithm\\
		Output: Sure, here is a Python implementation of the Fast Fourier Transform (FFT) algorithm:
		\begin{lstlisting}[language=python]
			import numpy as np
			
			def fft(x):
			n = len(x)
			if n == 1:
			return x
			even = fft(x[0::2])
			odd = fft(x[1::2])
			T = [np.exp(-2j * np.pi * k / n) * odd[k] for k in range(n // 2)]
			return [even[k] + T[k] for k in range(n // 2)] + [even[k] - T[k] for k in range(n // 2)]
		\end{lstlisting}
		The FFT is a widely used algorithm for computing the \textbf{Discrete Fourier Transform (DFT)} of a sequence. It reduces the number of computations needed from O(n$^2$) to O(n log n) 12. The basic idea of the FFT is to apply divide and conquer. We divide the coefficient vector of the polynomial into two vectors, recursively compute the DFT for each of them, and combine the results to compute the DFT of the complete polynomial 3.
		
		I hope this helps!
		
	\end{ReportAiUse}
	
\end{document}
%% 
%% This work consists of these files mcmthesis.dtx,
%%                                   figures/ and
%%                                   code/,
%% and the derived files             mcmthesis.cls,
%%                                   mcmthesis-demo.tex,
%%                                   README,
%%                                   LICENSE,
%%                                   mcmthesis.pdf and
%%                                   mcmthesis-demo.pdf.
%%
%% End of file `mcmthesis-demo.tex'.
